\section{Discussion}

The statistical characterization in \cref{obj:thm:fixed-binary-minimax} concerns the stability of a stationary reward functional across contracting four-state realizations. For fixed positive mixing scale and fixed nonnegative action-overlap exponent, the minimax mean-squared error has order \(T^{-1}\), where \(T\) is the trajectory length. The uniform modulus in \cref{obj:lem:pair-polynomial-modulus} explains how this rate accommodates latent-model degeneracies: agreement of seven scalar moments controls agreement of the stationary rewards, including at realizations with repeated transient eigenvalues or reduced effective dimension. The relevant stability is therefore a property of the stationary functional over the contracted class.

The explicit upper bound is \(C_{\alpha,L}/T\), where \(C_{\alpha,L}=B_\alpha^2(112V_{\alpha,L}+2)\). Since squared loss is at most one for the bounded procedures and targets used here, this displayed guarantee improves on the unit bound once \(T>C_{\alpha,L}\). The threshold can be large: \(B_\alpha\) grows as contraction slows and \(V_{\alpha,L}\) grows with the seventh power of the action likelihood-ratio bound. Thus the parametric exponent is a fixed-parameter statement; finite-sample usefulness also depends materially on mixing and action overlap.

Fixed joint-state dimension supplies the bounded recurrence underlying this conclusion. With four joint states, the transient characteristic polynomial has degree three, and comparing two realizations gives a transient polynomial of degree six. Contraction bounds these transient roots away from the stationary root at one. As expressed by \cref{obj:lem:pair-polynomial-modulus}, the resulting stationary-value modulus depends on the common contraction bound and remains uniform across the realizations. This connects the dimension of the joint-state space to the fixed number of moments used in \cref{obj:def:stable-grid-estimator}.

Action weighting gives this finite-dimensional structure an observed-data interpretation. The intervention identity in \cref{obj:lem:observable-intervention-moments} identifies the weighted reward moments with successive target-policy transitions starting from the behavior stationary law. Within each weighting window, the current action ratio supplies the target-policy reward vector, while the preceding ratios supply the target-policy transitions. Sequential assignment and action overlap justify this connection. For the empirical moments in \cref{obj:def:observable-moments}, the likelihood-ratio bound controls the weights and behavior contraction controls dependence across separated windows. The risk bound in \cref{obj:thm:stable-grid-upper} combines these sampling controls with stable extrapolation of the stationary functional.

The displayed estimator implements this characterization through a finite list of contracting moment fits. Its grid resolution \(M\) grows linearly with the trajectory length at a fixed mixing scale. The bounds in \cref{obj:prop:exhaustive-grid-arithmetic-complexity} quantify the size of the list before and after the contraction filter.

\begin{propositionv}{prop:exhaustive-grid-arithmetic-complexity}[Exact grid candidate counts]
Fix \(t_0>0\), the mixing-scale parameter, and define the associated contraction bound by
\[
\alpha=\exp(-1/t_0).
\]
For every integer \(M\ge0\), the four-coordinate simplex grid, represented by integer coordinates in \(\{0,\ldots,M\}\) summing to \(M\), has cardinality
\[
\binom{M+3}{3}.
\]
The unfiltered candidate list consists of an initial simplex vector \(\nu\), four simplex rows forming the transition candidate \(R\), and a reward vector \(r\) with four independently chosen grid coordinates. Its cardinality is exactly
\[
\binom{M+3}{3}^{5}(M+1)^4.
\]
For every integer \(M\ge0\) and every real contraction threshold \(\alpha\), retaining the candidates satisfying
\[
\delta(R)\le \alpha+\frac{6}{M},
\]
where \(\delta(R)\) is the Dobrushin contraction coefficient and division by zero is interpreted as zero, yields a cardinality at most
\[
\binom{M+3}{3}^{5}(M+1)^4.
\]

There exists a constant \(c>0\) such that, for every integer trajectory length \(T\ge1\), at the resolution
\[
M=\left\lceil(22+36/\alpha)T\right\rceil,
\qquad \alpha=\exp(-1/t_0),
\]
the unfiltered candidate-list cardinality and the filtered candidate-list cardinality are each at most \(cT^{19}\).
\end{propositionv}

The five simplex factors in \cref{obj:prop:exhaustive-grid-arithmetic-complexity} correspond to the initial probability vector and the four transition rows; the remaining factor counts the four reward coordinates. Each simplex contributes a cubic factor in the resolution, giving total degree nineteen. Thus candidate-list cardinality is the established computational descriptor of the displayed estimator. The exact product records the scale of exhaustive enumeration, and the \(cT^{19}\) bound expresses its growth with the horizon for a fixed mixing scale.

The statistical guarantee also applies within stationary-overlap subclasses of the binary experiment class. For any fixed stationary occupancy-ratio bound \(C>1\), impose the additional restriction \(d_e(s)\leq C d_b(s)\) at every joint state. The upper bound in \cref{obj:thm:stable-grid-upper} continues to apply uniformly on this restricted class. The two testing experiments in \cref{obj:thm:fixed-binary-minimax} satisfy the restriction because their stationary laws coincide, so they retain the matching lower bound. Consequently, for fixed positive mixing scale, fixed nonnegative action-overlap exponent, and \(T\geq12\), these subclasses also have minimax mean-squared-error order \(T^{-1}\).

Coincident policies give a direct sampling interpretation of the same order. Under the conditions of \cref{obj:prop:coincident-policy-reduction}, stationary initialization and behavior contraction make the ordinary reward sample mean unbiased and give a mean-squared-error bound proportional to \(T^{-1}\) for every \(T\geq1\). Together with the coincident-policy testing family in \cref{obj:def:four-state-testing-family}, this reduction connects the general stationary evaluation result to the familiar uncertainty of estimating a stationary reward mean.
